Knowledge editing, unlearning and belief implantation all assume that one targeted change can be made to a language model without collateral change.
Editing studies test this almost only on looked-up entity facts.
We ask how isolated a single \emph{computed} fact can be made: can a model answer 5 to \prompt{2+2=} and change nothing else?
On \llama we apply five counterfactual arithmetic edits with eleven edit methods, from a gated activation override to synthetic-document fine-tuning, and score each on about 990 must-not-change probes together with held-out paraphrases, entailments and True/False belief probes.
Exactly zero collateral change is reached only by the gated override, a keyed exception.
Its key need not be the literal string: a gate on a layer-20 activation covers 76\% of held-out paraphrases with zero flipped probes when the radius is chosen with hindsight, and 87\% with 0.8\% flips when it is chosen on held-out anchors.
No weight edit reaches zero.
The best ones, LoRA with a drift penalty, still flip 1.2--1.7\% of must-not-change probes, and the residue is systematic: probes whose true answer equals the edited fact's \emph{old} answer flip in 14--26\% of cases, against 0.5--1.3\% for all others.
Paraphrase generalization is cheap (94--97\%), but True/False belief stays at 1--10\%; only an in-context instruction and synthetic documents move it, and both leak more.
A reduced replication on \qwen agrees on the fine-tuning arms but not on the rank-one edit.
Our results imply that a locality score is hard to interpret unless the probe set contains facts that share the answer with the edited fact.
